\section*{Chapitre \ref{ch:rs:signal-combination} --- Combinaison de signaux}

\begin{solution}{exo:rs:signal-combination:1}
$10 \times 0.02/\sqrt{10 + 90 \times 0.25} = 0.2/\sqrt{32.5} = 0.035$~; la limite est $0.02/\sqrt{0.25} = 0.04$.
\end{solution}

\begin{solution}{exo:rs:signal-combination:2}
Les poids pondérés par l'IC utilisent les parties positives 0.03, 0.01 et 0~: poids 0.75, 0.25 et 0. À mi-chemin vers des poids équipondérés~:
$\frac12(0.75, 0.25, 0) + \frac12(\frac13, \frac13, \frac13) = (0.542, 0.292, 0.167)$.
\end{solution}

\begin{solution}{exo:rs:signal-combination:3}
Les rendements et les signaux sont ordonnés dans le temps, corrélés entre dates voisines et exposés à des régimes~; un pli qui apprend sur des dates postérieures et teste sur
des dates antérieures permet aux poids de \omterm{def:rs:signal-combination:stacking}{stacking} d'apprendre du futur (quel mélange a fonctionné dans le régime propre à la période de test). Les fenêtres croissantes qui prédisent toujours
vers l'avant reproduisent ce que la firme peut faire en production.
\end{solution}

\begin{solution}{exo:rs:signal-combination:4}
$w \propto C^{-1}\rho$ avec $C = \begin{pmatrix} 1 & 0.5\\ 0.5 & 1\end{pmatrix}$~: $C^{-1}\rho = \frac{1}{0.75}(0.03 - 0.01, -0.015 + 0.02) = (0.0267, 0.0067)$,
poids 0.8 et 0.2. Le meilleur IC est $\sqrt{\rho^\top C^{-1}\rho} = \sqrt{0.03 \times 0.0267 + 0.02 \times 0.0067} = 0.0306$~; des poids équipondérés donnent
$0.05/\sqrt{2 + 2 \times 0.5} = 0.0289$.
\end{solution}

\begin{solution}{exo:rs:signal-combination:5}
Le premier signal est son propre premier vecteur de Gram--Schmidt. Le second, résidualisé sur le premier, est $z_2 - cz_1$, de variance $1 - c^2$~; standardisé,
$(z_2 - cz_1)/\sqrt{1 - c^2}$. Sa corrélation avec la cible est $(\rho_2 - c\rho_1)/\sqrt{1 - c^2}$~: ce que le second signal sait et que le premier
ne sait pas.
\end{solution}

\begin{solution}{exo:rs:signal-combination:6}
Soixante observations de quarante séries donnent une matrice de covariance empirique dont les plus petites valeurs propres sont très inférieures aux vraies (la dispersion des valeurs propres empiriques
du livre~4, chapitre~22), et $\Sigma^{-1}\mu$ divise par elles~: des directions qui paraissent presque sans risque à cause de l'erreur d'échantillonnage reçoivent les
poids les plus grands. Le rétrécissement vers la diagonale relève ces valeurs propres~; à la diagonale, la règle devient la moyenne sur la variance par signal, proche des poids pondérés par l'IC.
\end{solution}

\begin{solution}{exo:rs:signal-combination:7}
Le $\lambda$ optimal (sur une grille de dixièmes) vaut 0.5 à 6 mois (IC 0.074), 0.4 à 12, 0.6 à 24 et 0.4 à 48 (0.083). Il ne se déplace pas régulièrement vers
les poids de régression~: avec 100 titres, même quatre années laissent assez d'erreur dans quarante poids de régression pour qu'un fort rappel vers des poids équipondérés rapporte, et
l'optimum est bruité d'une fenêtre à la suivante.
\end{solution}

\begin{solution}{exo:rs:signal-combination:8}
Les poids ont été ajustés sur la période sur laquelle on les juge~: la comparaison mesure l'ajustement, non la prévision. Réajuster sur des dates antérieures à chaque date jugée
(fenêtres croissantes ou plis ordonnés dans le temps), comparer hors échantillon (out-of-sample), et présenter aussi les poids équipondérés~; sur le panel complet du chapitre, l'avantage dans l'échantillon (in-sample)
du \omterm{def:rs:signal-combination:weights}{mélange à ICIR maximal} non rétréci s'est inversé hors échantillon (out-of-sample) dans le monde en dérive.
\end{solution}

\begin{solution}{pb:rs:signal-combination:1}
\begin{enumerate}
\item 0.015 et 0.027.
\item 0.065 d'après la proposition~; 0.064 mesuré.
\item $\bar\rho/\sqrt{\bar c} = 0.090$. Chaque nouveau signal reprend une partie du bruit de son thème, partagé avec les autres signaux du thème~; cette part commune ne
s'annule pas en moyenne, et l'IC du mélange n'approche le plafond que lorsque le nombre de signaux croît sans borne.
\item Cinq signaux par thème partagent le bruit du thème.
\item Monde stable~: 0.064, 0.099 et 0.102. Monde en dérive~: 0.060, 0.075 et 0.087.
\item Il inverse la matrice de covariance des quarante séries d'IC estimées sur 59 mois et charge sur ses plus petites valeurs propres, les plus sous-estimées.
\item 0.098 dans le monde stable et 0.074 dans le monde en dérive~: proche du résultat pondéré par l'IC.
\item L'information des thèmes dérive avec une demi-vie d'environ deux ans, si bien que les cinq premières années disent encore quelque chose des suivantes, et 704 titres
estiment les poids avec précision.
\item 0.073 dans le monde stable, 0.005 dans le monde en dérive.
\item \textbf{Résultat nommé.} Avec 100 titres et un an d'historique, le meilleur mélange hors échantillon (out-of-sample) dans le monde stable est la pondération par l'IC (0.086 contre 0.069
pour les poids équipondérés), et les poids de régression sont mieux rétrécis de 40~\% vers des poids équipondérés (0.081)~; dans le monde en dérive, le meilleur mélange est l'équipondération
($\lambda = 1$, 0.069), et les poids de régression non rétrécis rapportent 0.005.
\item Monde stable~: 0.085, 0.086, 0.082 et 0.092 contre 0.069, meilleur à chaque fenêtre. Monde en dérive~: 0.043, 0.049, 0.067 et 0.074, moins bon jusqu'à quatre
années d'historique.
\item L'erreur d'estimation à taille d'échantillon finie des poids de combinaison.
\item Tout son poids sur la pondération par l'IC, dans les deux mondes. Il choisit parmi ses mélanges de base sur des dates retenues~; il ne peut pas être meilleur que le meilleur d'entre eux de plus
que ne le permet la diversification de leurs erreurs, et ici une seule base dominait.
\item 0.064 pour l'ensemble symétrique, comme pour les signaux bruts~; 0.050 pour l'ensemble séquentiel. L'ordre séquentiel donne au meilleur signal la
part commune de son thème et laisse aux autres un bruit résiduel, que les poids équipondérés surpondèrent ensuite~; l'ensemble symétrique modifie chaque signal le moins possible.
\item Quand les poids doivent être interprétés (information marginale), quand les expositions du portefeuille doivent être retirées (résidualisation sur le modèle de risque), ou quand un nouveau
signal doit être jugé sur ce qu'il apporte~; pas pour créer de l'information.
\item Sur des plis ordonnés dans le temps à l'intérieur de la période d'apprentissage, en choisissant le $\lambda$ qui maximise l'IC moyen hors pli, réajusté selon un calendrier, avec
les poids équipondérés comme valeur par défaut lorsque le choix n'est pas nettement meilleur.
\item Des signaux de qualité très différente et persistante, ou des signaux à IC négatifs~: les poids équipondérés les conservent tous.
\item Le panel est grand (704 titres par mois), si bien que la régression non pénalisée est déjà précise, et pénaliser vers zéro ne fait surtout que remettre à l'échelle un composite
qui est de toute façon standardisé.
\item L'IC du \omterm{def:rs:signal-combination:blend}{mélange équipondéré} sur les mêmes dates, dans l'échantillon (in-sample) et hors échantillon (out-of-sample).
\item Quand la qualité des signaux diffère de façon persistante et que l'échantillon est assez grand pour que l'erreur d'estimation des poids soit plus petite que les écarts
qu'ils estiment.
\end{enumerate}
\end{solution}

\begin{solution}{iq:rs:signal-combination:1}
Standardiser chacun par date, partir des poids équipondérés et mesurer l'IC du mélange hors échantillon (out-of-sample)~; essayer ensuite les poids pondérés par l'IC et les poids de régression rétrécis vers
des poids équipondérés, le rétrécissement étant choisi sur des plis ordonnés dans le temps~; ne garder les poids estimés que s'ils battent les poids équipondérés hors échantillon (out-of-sample) de plus que
leur bruit. Vérifier d'abord la structure de corrélation des signaux~: c'est elle qui fixe le plafond du mélange.
\end{solution}

\begin{solution}{iq:rs:signal-combination:2}
Les moyennes simples de prévisions battent à plusieurs reprises les combinaisons à poids estimés (Stock et Watson~; Smith et Wallis). L'explication~: l'erreur dans
l'estimation des poids à partir d'échantillons finis, et l'instabilité de la qualité relative des prévisions, peut dépasser le gain de la pondération.
\end{solution}

\begin{solution}{iq:rs:signal-combination:3}
Construire les mélanges de base sur des fenêtres croissantes, prédire avec chacun le bloc de dates suivant, et ajuster des poids de \omterm{def:rs:signal-combination:stacking}{stacking} non négatifs
uniquement sur les prévisions hors échantillon (out-of-sample)~; ne jamais laisser une base ou le stacker voir des dates sur lesquelles il sera jugé plus tard, et comparer le stack à sa meilleure base et aux poids équipondérés.
\end{solution}

\begin{solution}{iq:rs:signal-combination:4}
Il peut répéter une information déjà présente dans le composite (forte corrélation avec les signaux existants~: son IC résiduel, $(\rho - c\rho_1)/\sqrt{1 - c^2}$ pour un
signal corrélé, peut être proche de zéro)~; il peut être plus bruité que ne le suggère son IC~; ou la pondération peut lui avoir donné trop peu de poids pour compter. Mesurer son
IC incrémental après résidualisation sur le composite (chapitre~12).
\end{solution}

\begin{solution}{iq:rs:signal-combination:5}
L'\omterm{def:rs:signal-combination:orth}{orthogonalisation} séquentielle (Gram--Schmidt) garde le premier signal et résidualise chaque signal suivant sur ceux qui le précèdent~: l'ordre décide qui garde
l'information commune, utile lorsqu'il existe une priorité naturelle (les facteurs d'un portefeuille d'abord, les nouveaux signaux ensuite). L'\omterm{def:rs:signal-combination:orth}{orthogonalisation symétrique} ($ZS^{-1/2}$) traite tous les
signaux de la même manière et donne l'ensemble décorrélé le plus proche des originaux~: utile pour l'interprétation sans ordre.
\end{solution}

\begin{solution}{iq:rs:signal-combination:6}
Maximiser $w^\top\rho/\sqrt{w^\top Cw}$~: la condition de gradient donne $\rho(w^\top Cw) = Cw(w^\top\rho)$, donc $w \propto C^{-1}\rho$ et le maximum vaut
$\sqrt{\rho^\top C^{-1}\rho}$. L'inverse estimé $C^{-1}$ amplifie l'erreur d'échantillonnage des plus petites valeurs propres de $C$ et $\hat\rho$ est lui-même bruité, si bien que les
poids sont extrêmes et instables~; rétrécir $C$ vers sa diagonale et $w$ vers des poids équipondérés.
\end{solution}
